\section{Apertura}

% ── Frame: Contexto ───────────────────────────────────────────────────────────
\begin{frame}{¿Quién y por qué?}
  \begin{columns}[T]
    \begin{column}{0.58\textwidth}
      {\color{SotmVerde}\bfseries ¿Quién?}
      \begin{itemize}
        \item Hernán Galileo Cabrera Garibaldi.
        % \pause
        \item Ingeniero en Computación y Maestrante de Urbanismo; UNAM.
        \item Fundador y Desarrollador de Apimetro.   % EDITAR: descomentar si aplica
        % \pause
      \end{itemize}

      \vspace{0.2cm}
      {\color{SotmVerde}\bfseries ¿Por qué?}
      \begin{itemize}
        % EDITAR: borradores tomados del artículo; ajusta la redacción
        \item Uno de los grandes desafíos en el \textbf{World Building} es la creación de un mundo congruente para sus personajes y sus necesidades.
        % \pause
        \item Abrir el dato no es democratizarlo.
      \end{itemize}
    \end{column}
    \begin{column}{0.38\textwidth}
      \begin{block}{Pregunta Gancho}
        ¿Cuántas fuentes de datos necesitas para trazar un viaje de Ecatepec (EDOMEX) a Milpa Alta (CDMX)?
      \end{block}
    \end{column}
  \end{columns}
\end{frame}

% ── Frame: Respuesta al gancho — flujo manual sin Apimetro ───────────────────
\begin{frame}{Respuesta a la Pregunta Gancho}
  \centering
  \begin{tikzpicture}[
    fuente/.style   = {rectangle, rounded corners=5pt,
                       draw=SotmVerde, line width=0.8pt,
                       fill=SotmCrema, text width=3.8cm,
                       align=left, minimum height=0.85cm, font=\small},
    num/.style      = {circle, fill=SotmVerde, text=white,
                       font=\small\bfseries, inner sep=0pt, minimum size=0.5cm},
    caja/.style     = {rectangle, rounded corners=5pt,
                       draw=SotmVerde, line width=0.8pt,
                       fill=SotmCrema, text width=2.4cm,
                       align=center, minimum height=0.8cm, font=\small},
    result/.style   = {caja, fill=SotmTerracota, draw=SotmTerracota,
                       text=white, font=\small\bfseries},
    decision/.style = {diamond, draw=SotmVerde, line width=0.8pt,
                       fill=SotmCrema, aspect=2, inner sep=1pt,
                       text width=2.0cm, align=center, font=\scriptsize},
    flecha/.style   = {-Stealth, thick, color=SotmVerde},
    lateral/.style  = {-Stealth, thick, color=SotmTerracota, dashed}
  ]
    % ── Fuentes: la respuesta a "¿cuántas?" ──
    \node[fuente] (f1) at (0,  1.05) {\textbf{GTFS SEMOVI}\\{\scriptsize Metro · Metrobús · RTP}};
    \node[fuente] (f2) at (0,  0)    {\textbf{OSM}\\{\scriptsize Mexibús · Mexicable}};
    \node[fuente] (f3) at (0, -1.05) {\textbf{Trazos a mano}\\{\scriptsize Trolebús L13 · colectivos}};
    \foreach \i in {1,2,3}
      \node[num, left=0.15cm of f\i] {\i};

    % ── Validación y resultado ──
    \node[decision] (dec)  at (5.0, 0)    {¿Los trazos coinciden?};
    \node[result]   (res)  at (9.1, 0)    {Iniciar el viaje};
    \node[caja]     (corr) at (5.0, -1.95) {Corrección};

    \foreach \i in {1,2,3}
      \draw[flecha] (f\i.east) -- (dec.west);
    \draw[flecha]  (dec)  -- node[above, font=\scriptsize]{Sí} (res);
    \draw[lateral] (dec)  -- node[right, font=\scriptsize]{No} (corr);
    \draw[lateral] (corr) -| (f3.south);
  \end{tikzpicture}

  \vspace{0.1cm}
  {\color{SotmTerracota}\rule{0.6\textwidth}{0.8pt}}\\[0.15cm]
  {\bfseries\color{SotmVerde} $\geq 3$ fuentes · formatos incompatibles · trabajo manual}\\[0.1cm]
  {\small\itshape …¿y si fuera una sola consulta?}
\end{frame}
